\section{Digression sur le concept de sélection causale dans le domaine des sciences cognitives}\label{scienceco}

De nombreuses recherches en sciences cognitives ont été menées pour tenter de modéliser la façon dont nous concevons et préférons certains scénarios contrefactuels par rapport à d'autres. Ces mécanismes sont influencés par différents critères tels que la structure causale, le temps, les normes, la valeur de l'effet produit... Même si certaines études portent, plus ou moins implicitement, sur la notion de responsabilité causale plus que sur le concept de sélection causale, des définitions et concepts peuvent être utiles pour définir et orienter les critères d'évaluation ou encore de construction de scénarios contrefactuels. Je présente ci-dessous une liste non exhaustive de notions trouvées dans la littérature pouvant être potentiellement intégrées à notre cadre.

\subsection{Mutabilité et propensité des facteurs}\label{mut_prop}
Les notions de mutabilité (à quel point il est simple d'imaginer qu'un facteur ait pu être différent) et de propensité (la vraisemblance qu'un effet ait lieu sachant une cause) jouent, selon \textcite{kominsky_causal_2015} et \textcite{mcgill_mutability_2000}, un rôle essentiel dans le raisonnement contrefactuel. Le premier article vise à mettre en évidence, dans le processus d'évaluation causale d'un événement A, l'importance de la mutabilité et de la propensité d'un événement B nécessaire à la production de l'effet (structure causale conjointe). D'autre part dans \textcite{kominsky_causal_2015} le travail porte cette fois-ci sur le phénomène de \textit{"causal superseding"} qui est décrit comme : le phénomène selon lequel la violation d'une norme par un agent diminue la causalité attribuée à un autre agent. Même si cet effet sort de notre cadre par son aspect multiagents, on trouve des résultats similaires à ceux de \textcite{mcgill_mutability_2000}. Ils avancent, eux aussi, que selon leur modèle, un facteur causal est vu comme plus faible lorsqu'un autre facteur est plus mutable, mais uniquement si ce facteur est très susceptible de produire l'effet. La mutabilité a été introduite dans la fonction d'évaluation proposée dans la section \ref{fct_cout}.

\subsection{Impact des normes sur le raisonnement contrefactuel}
Dans \parencite{icard_normality_2017} une mesure de force de causalité (correspondant à un compromis entre suffisance et nécessité) est présentée et permet de modéliser les effets de la normalité sur le jugement causal. Dans ce travail aucune distinction de types de normes (statistique ou prescriptive) n'est appliquée et les auteurs travaillent dans un cadre non déterministe ne faisant pas de différence entre les états et les événements. Cette mesure n'est donc pas directement applicable à notre cadre, mais les effets discutés présentent des intérêts potentiels. Une différence des effets de la normalité est faite selon la structure causale. Dans les cas conjonctifs, on observe un effet d'\textit{abnormal inflation} décrit par le fait qu'une cause soit considérée comme plus causale si elle est anormale par rapport au cas où elle est normale. Le second effet pour les structures conjonctives est l'effet de \textit{supersession}: un facteur est considéré comme plus causal lorsque les autres facteurs sont normaux. Dans le cas disjonctif, on observe l'effet d'\textit{abnormal deflation}: une cause est considérée comme plus causale lorsque qu'elle est normale que lorsqu'elle ne l'est pas. Et enfin toujours dans le cas disjonctif l'effet de \textit{no supersession with disjunction} qui avance qu'une cause  est tout aussi encline à être jugée comme causale lorsqu'une seconde cause est normale, que lorsqu'elle est anormale.

\subsection{Impact du temps sur le raisonnement contrefactuel}

Dans \parencite{henne_counterfactual_2021}, le travail porte sur l'influence qu'a la position temporelle d'un événement dans une séquence causale sur le jugement de sa responsabilité, et comment le raisonnement contrefactuel peut expliquer ces effets. Le premier effet discuté est l'effet de récence (dans un match, on juge plutôt le dernier but comme la cause de la victoire plutôt que le premier). Mais les auteurs observent que dans des cas de surdétermination et plus particulièrement de préemption tardive, cet effet s'inverse : si chaque effet est suffisant alors la causalité est attribuée à l'événement le plus ancien. Ils étendent le \textit{necessity-sufficiency model} présenté par \textcite{icard_normality_2017}, qui avance que plus un événement fait une différence dans une simulation contrefactuelle, plus il est vu comme causal, en introduisant l'hypothèse selon laquelle les gens sont plus enclins à imaginer des contrefactuels autour de l'événement le plus récent (quelle que soit la structure). Avec cette extension du modèle causal, le mécanisme est toujours le même (simuler le contrefactuel de l'événement le plus récent) mais c'est la structure causale qui détermine ce que la pensée contrefactuelle révèle. Dans un cas, elle indique que l'événement causal est le plus récent et dans l'autre c'est le plus ancien.

